\documentclass[11pt]{amsart}
\usepackage[margin=1.1in]{geometry}
\usepackage{amsmath,amssymb,amsthm}
\usepackage{booktabs}
\usepackage{enumitem}
\usepackage[hidelinks]{hyperref}

\newtheorem{theorem}{Theorem}[section]
\newtheorem{lemma}[theorem]{Lemma}
\newtheorem{proposition}[theorem]{Proposition}
\newtheorem{corollary}[theorem]{Corollary}
\theoremstyle{definition}
\newtheorem{definition}[theorem]{Definition}
\newtheorem{example}[theorem]{Example}
\theoremstyle{remark}
\newtheorem{remark}[theorem]{Remark}

\newcommand{\F}{\mathbb{F}}
\newcommand{\Aut}{\operatorname{Aut}}
\newcommand{\gon}{\operatorname{gon}}
\newcommand{\Pl}{\mathcal{P}}
\newcommand{\LL}{\mathcal{L}}
\newcommand{\minpoly}{\operatorname{mp}}
\newcommand{\Enc}{\operatorname{Enc}}
\newcommand{\code}[1]{\texttt{#1}}

\title[Orbits on places of small degree]{Orbits of $\langle\sigma_4\rangle$ on places of small degree of $X_1(42)_{\F_5}$:\\ mathematics and computational details}
\author{Sourajyoti Maiti}
\address{Department of Mathematics, University of Connecticut}
\email{sourajyoti.maiti@uconn.edu}
\date{\today}

\begin{document}
\sloppy\emergencystretch=3em

\begin{abstract}
This note documents the computation of the orbits of the cyclic group $G=\langle\sigma_4\rangle$ of order $12$ of automorphisms of $X=X_1(42)_{\F_5}$ on the places of degree $1,\dots,5$, and explains how the result is stored, labelled independently of the Magma session, and used in the orbit-reduced search of the manuscript ``The $\mathbb{Q}$-gonality of $X_1(42)$ is $12$'' (Section~6 there). Every step is exact. The image of a place is found from the pull-back $\sigma^*$ alone, with no need to compute the inverse automorphism or any point map. The note is a companion to the scripts \code{X1\_42\_orbits.m}, \code{X1\_42\_orbits\_merge.py}, \code{run\_orbits.sh}, \code{X1\_42\_orbits\_load.m}, \code{X1\_42\_orbits\_check.m}, \code{X1\_42\_rational\_labels.m} and \code{X1\_42\_split.py}; it is not part of the manuscript.
\end{abstract}

\maketitle
\tableofcontents

%======================================================================
\section{Setup}
%======================================================================

\subsection{The curve}
Let $X=X_1(42)_{\F_5}$; it has genus $25$. We use the reduction modulo $5$ of Sutherland's plane model $f(x,y)=0$ (file \code{FFFc42.txt}); see Lemma~2.5 of the manuscript. Magma builds the function field as
\[
\code{AFF := AlgorithmicFunctionField(FunctionField(XF2))},\qquad
K=\F_5(X)=\F_5(t)[w]/(\Phi(t,w)),\qquad [K:\F_5(t)]=16,
\]
with $t=$ \code{BaseField(AFF).1} and $w=$ \code{AFF.1}. Here $XF2$ is the projective closure of $\bar f(x,y)=0$, and Magma works in the affine chart $y\neq0$, so that
\[
t=\frac1y,\qquad w=\frac xy,\qquad \Phi(t,w)=t^{23}\,\bar f\Bigl(\frac wt,\frac1t\Bigr),
\]
as in Appendix~A.1 of the manuscript (checked by \code{X1\_42\_check\_tw3.m}, which also writes $\Phi$ to \code{X1\_42\_Phi.txt}).

\begin{remark}[Notation in the scripts]
The scripts call the generator $w$ ``\code{y}'' (variables \code{YY}, \code{SL\_YY}, \code{GY}), and $t$ is \code{TT} or \code{SL\_TT}. Throughout this note we use the notation $t,w$ of the manuscript. The model coordinates $x,y$ do not appear below.
\end{remark}

All places are places of $K/\F_5$. For $v\ge1$ let $\Pl_v$ be the list \code{Places(AFF,v)} of places of degree $v$, in the order in which Magma returns it, and write $P^{(v)}_i$ for its $i$-th entry. This order is \emph{not} the same in every Magma session (Section~6.3 of the manuscript); Sections~\ref{sec:checks} and~\ref{sec:labels} explain how this is handled. The sizes are
\[
\begin{array}{c|ccccccccc}
v & 1&2&3&4&5&6&7&8&9\\\hline
\#\Pl_v & 12&6&8&183&648&2288&11160&48708&214384
\end{array}
\]
The present computation covers $v\le5$.

\subsection{The group}
By Proposition~A.1 of the manuscript, $\sigma=\sigma_4$ generates a cyclic group $G$ of order $12$; the manuscript uses only $G$ (Magma reports that $G$ is the whole of $\Aut_{\F_5}(X)$, but this is not used). The automorphisms are stored in \code{X1\_42\_aut\_F5.txt} as the pairs $(\sigma_i^*t,\sigma_i^*w)\in K^2$ of pull-backs. The scripts never recompute them: they read the file with \code{LoadAutDataOpt(name,false)}, which takes about $2$\,s and skips the closure test. They then check that $\sigma_4$ has order exactly $12$, by composing $\sigma_4^*$ with itself until the identity $(t,w)$ appears; this takes about $87$\,s.

As in Section~4 of the manuscript, Magma returns field automorphisms; for a field automorphism $\sigma$ we write $\varphi$ for the automorphism of $X$ with $\varphi^*=\sigma$, and $\sigma(P)$ means $\varphi(P)$.

On the rational places, $\sigma_4$ acts as a $12$-cycle. In the session-independent labels of Section~\ref{sec:labels} (file \code{X1\_42\_rational\_labels.txt}) it is
\begin{equation}\label{eq:S}
S=[12,9,7,3,6,10,1,2,5,11,4,8]=(1\,12\,8\,2\,9\,5\,6\,10\,11\,4\,3\,7).
\end{equation}
In the numbering of \code{Places(AFF,1)} in the session of \code{X1\_42\_check\_orbs.m} it was recorded as
\begin{equation}\label{eq:A4}
A_4=[9,5,2,1,12,8,11,7,3,4,10,6]=(1\,9\,3\,2\,5\,12\,6\,8\,7\,11\,10\,4).
\end{equation}
Only \eqref{eq:S} is used by the search; \eqref{eq:A4} serves as a consistency check (Section~\ref{sec:checks}).

%======================================================================
\section{Automorphisms acting on places}
%======================================================================

\subsection{From the pull-back to places}
Let $\phi\in G$ and write $\phi^*\colon K\to K$ for the corresponding $\F_5$-automorphism of the function field. For a place $P$ with valuation ring $\mathcal{O}_P$, put
\[
\mathcal{O}_{\phi(P)}=\{f\in K:\ \phi^*f\in\mathcal{O}_P\}=(\phi^*)^{-1}(\mathcal{O}_P).
\]
Since $\phi^*$ is a field automorphism fixing $\F_5$, this is again a valuation ring of $K/\F_5$. It defines a place $\phi(P)$.

\begin{lemma}\label{lem:basic}
For every place $P$ and every $f\in K$:
\begin{enumerate}[label=\textup{(\alph*)}]
\item $v_{\phi(P)}(f)=v_P(\phi^*f)$;
\item $\phi^*$ induces an $\F_5$-isomorphism $\bar\phi\colon\kappa(\phi(P))\to\kappa(P)$ of residue fields, with $\bar\phi(f(\phi(P)))=(\phi^*f)(P)$ whenever $f\in\mathcal{O}_{\phi(P)}$. In particular $\deg\phi(P)=\deg P$;
\item $P\mapsto\phi(P)$ is a bijection on the places of each degree $v$, and $\phi\mapsto(P\mapsto\phi(P))$ is an action of $G$, possibly through the opposite group, which does not change the orbits.
\end{enumerate}
\end{lemma}

\begin{proof}
(a) holds because $v_P\circ\phi^*$ is a normalised discrete valuation of $K$ with valuation ring $\mathcal{O}_{\phi(P)}$. For (b), $\phi^*$ maps $\mathcal{O}_{\phi(P)}$ onto $\mathcal{O}_P$ and, by (a), maps the maximal ideal onto the maximal ideal. It fixes $\F_5$, so the induced map of residue fields is an $\F_5$-isomorphism, and the degree is preserved. For (c), $(\phi^{-1})^*$ provides the inverse map, and $(\phi\psi)^*=\psi^*\phi^*$ shows that composition is respected up to order.
\end{proof}

Because $G=\langle\sigma\rangle$ is cyclic, the $G$-orbit of a place $P$ is $\{\sigma^k(P):0\le k<12\}$. Let $\pi_v$ be the permutation of $\{1,\dots,\#\Pl_v\}$ with $\sigma(P^{(v)}_i)=P^{(v)}_{\pi_v(i)}$. Then:

\begin{corollary}\label{cor:cycles}
The $G$-orbits on $\Pl_v$ are exactly the cycles of $\pi_v$. Every cycle length divides $12$, and $\pi_v^{12}=\mathrm{id}$.
\end{corollary}

So it suffices to compute the single permutation $\pi_v$ for each degree $v$. The difficulty is that $\sigma$ is known only through $\sigma^*$. Computing $\sigma(P)$ directly as a place would require the ideal $(\sigma^*)^{-1}(\mathcal{O}_P)$. Instead, we \emph{identify} $\sigma(P)$ inside the already computed list $\Pl_v$, using only valuations and residues of explicit functions at $P$. By Lemma~\ref{lem:basic}, these are values at $\sigma(P)$ in disguise.

%======================================================================
\section{Identifying the image of a place}
%======================================================================

Fix a degree $v$ and a place $P\in\Pl_v$. The image $\sigma(P)$ is found in at most three exact steps. Each step shrinks a set of candidates that is \emph{guaranteed} to contain $\sigma(P)$.

\subsection{Encoding polynomials}
For a polynomial $g=\sum_{i=0}^{n}c_ix^i\in\F_5[x]$ with $c_n\neq0$, put
\[
\Enc(g)=\sum_{i=0}^{n}c_i5^i+5^{n+1}\in\mathbb{Z}_{>0},\qquad c_i\in\{0,\dots,4\}.
\]
$\Enc$ is injective: the leading power $5^{n+1}$ recovers $n$, and the base-$5$ digits below it recover the coefficients. For monic $g$ this is the encoding $\operatorname{enc}$ of Section~6.3 of the manuscript. For $\alpha$ in a finite extension of $\F_5$, let $\minpoly(\alpha)\in\F_5[x]$ denote its minimal polynomial over $\F_5$.

\subsection{Step 1: the bucket}
The place of $\F_5(t)$ under a place $Q$ of $K$ is determined by $\minpoly(t(Q))$ if $t$ is regular at $Q$, and is the infinite place if $t$ has a pole at $Q$. Define the \emph{bucket} of $Q$ as
\[
b(Q)=\begin{cases}\Enc(\minpoly(t(Q))) & v_Q(t)\ge0,\\ 0 & v_Q(t)<0.\end{cases}
\]
We compute $b(Q)$ for all $Q\in\Pl_v$ once, and group the indices by bucket.

\begin{lemma}\label{lem:bucket}
Put $b^\sigma(P)=\Enc(\minpoly((\sigma^*t)(P)))$ if $v_P(\sigma^*t)\ge0$, and $b^\sigma(P)=0$ otherwise. Then $b(\sigma(P))=b^\sigma(P)$.
\end{lemma}

\begin{proof}
By Lemma~\ref{lem:basic}(a), $v_{\sigma(P)}(t)=v_P(\sigma^*t)$, so $t$ has a pole at $\sigma(P)$ iff $\sigma^*t$ has a pole at $P$. Otherwise $\bar\sigma(t(\sigma(P)))=(\sigma^*t)(P)$ by Lemma~\ref{lem:basic}(b). An $\F_5$-isomorphism of fields preserves minimal polynomials over $\F_5$.
\end{proof}

So the first candidate set is $C_1=\{j:\ b(P^{(v)}_j)=b^\sigma(P)\}$. It contains $\pi_v(i)$ for $P=P^{(v)}_i$, and has at most $16$ elements, since $\sum e_if_i=[K:\F_5(t)]=16$ over any place of $\F_5(t)$. If $\#C_1=1$, the image is found.

\subsection{Step 2: the key}
For a place $Q$ at which $t$ and $w$ are both regular, put $a=t(Q)$ and $c=w(Q)$, and define its \emph{key}
\[
\kappa(Q)=\sum_{m=0}^{3}\Enc(g_m)\,(5^{12})^{3-m},\qquad (g_0,g_1,g_2,g_3)=\bigl(\minpoly(a),\minpoly(c),\minpoly(a+c),\minpoly(a+2c)\bigr).
\]
If $t$ or $w$ has a pole at $Q$, we set $\kappa(Q)=0$, meaning ``not available''. All $g_m$ have degree at most $v\le9$, so $\Enc(g_m)<5^{11}<5^{12}$, and $\kappa$ is an injective function of the quadruple $(g_0,\dots,g_3)$.

\begin{lemma}\label{lem:key}
If $\sigma^*t$ and $\sigma^*w$ are regular at $P$, then $\kappa(\sigma(P))$ equals the key computed from the pair $\bigl((\sigma^*t)(P),(\sigma^*w)(P)\bigr)$ in place of $(a,c)$.
\end{lemma}

\begin{proof}
Use the same isomorphism $\bar\sigma$ for both coordinates. It sends $(t(\sigma(P)),w(\sigma(P)))$ to $((\sigma^*t)(P),(\sigma^*w)(P))$, and hence sends $a+c$ and $a+2c$ to the corresponding combinations. Minimal polynomials over $\F_5$ are preserved.
\end{proof}

The second candidate set is $C_2=\{j\in C_1:\ \kappa(P^{(v)}_j)=\kappa^\sigma(P)\}$. It still contains the true image. If $\#C_2=1$, the image is found. If $C_2$ is empty (which can only happen when a key is unavailable), or if $\kappa^\sigma(P)$ cannot be computed, we continue with $C=C_1$; otherwise we continue with $C=C_2$.

\begin{remark}\label{rem:collisions}
Two different places can have the same pair $(b,\kappa)$: for example, two places over the same place of $\F_5(t)$ whose residues $(a,c)$ are conjugate over $\F_5$ only coordinatewise, but not as a pair. This happens for $10$ of the $12$ rational places. For the places of degree $2\le v\le5$ it does not happen: the pairs $(b(Q),\kappa(Q))$ are pairwise distinct in each of these degrees (checked by \code{X1\_42\_split.py} and \code{X1\_42\_rational\_labels.m}; Section~6.3 of the manuscript). Step~3 handles whatever remains.
\end{remark}

\subsection{Step 3: the uniformizer pattern}
Let $C=\{j_1,\dots,j_r\}$ with $r\ge2$, and write $Q_s=P^{(v)}_{j_s}$. For every candidate $Q_s$, choose a function $u_s\in K$ with $v_{Q_s}(u_s)>0$. In practice $u_s=\code{LocalUniformizer}(Q_s)$, and the script checks $v_{Q_s}(u_s)>0$. Nothing is required of $u_s$ at the other candidates. Record:
\begin{itemize}
\item the zero pattern $Z_s=\{q\in\{1,\dots,r\}:\ v_{Q_q}(u_s)>0\}$, which always contains $s$;
\item the test bit $\beta_s=\bigl[\,v_P(\sigma^*u_s)>0\,\bigr]\in\{0,1\}$.
\end{itemize}
A candidate $q$ \emph{survives} if $q\in Z_s\iff\beta_s=1$ for every $s=1,\dots,r$.

\begin{proposition}\label{prop:unif}
Let $q^\ast$ be the index with $Q_{q^\ast}=\sigma(P)$.
\begin{enumerate}[label=\textup{(\alph*)}]
\item $q^\ast$ always survives.
\item If exactly one candidate survives, it is $q^\ast$.
\item A candidate $q\neq q^\ast$ survives only if $u_q$ vanishes at $\sigma(P)$.
\end{enumerate}
\end{proposition}

\begin{proof}
By Lemma~\ref{lem:basic}(a), $\beta_s=[v_{\sigma(P)}(u_s)>0]=[q^\ast\in Z_s]$, which proves (a). (b) follows from (a). For (c): $q\in Z_q$, so $q$ can survive only if $\beta_q=1$, i.e. $v_{\sigma(P)}(u_q)>0$.
\end{proof}

So Step 3 is never wrong; at worst it is inconclusive. In that case the image is recorded as $-1$ (\emph{ambiguous}), and the whole degree is rejected at the merge stage (Section~\ref{sec:merge}). In the computation this did not occur: Steps~1--3 determined $\sigma(P)$ for every place of degree $\le5$.

\begin{example}
Suppose $C=\{1,2,3\}$, with zero patterns $Z_1=\{1\}$, $Z_2=\{2,3\}$ and $Z_3=\{3\}$, and observed bits $(\beta_1,\beta_2,\beta_3)=(0,1,1)$. Candidate $1$ fails at $s=1$, since $1\in Z_1$ but $\beta_1=0$. Candidate $2$ fails at $s=3$, since $2\notin Z_3$ but $\beta_3=1$. Candidate $3$ satisfies all three conditions. Hence $\sigma(P)=Q_3$. With bits $(0,1,0)$ instead, only candidate $2$ survives.
\end{example}

\subsection{Cost}
The functions $\sigma^*t$, $\sigma^*w$ and $\sigma^*u_s=u_s(\sigma^*t,\sigma^*w)$ (\code{ApplyAut}) are exact elements of $K$. Steps 1 and 2 need two residue evaluations and four minimal polynomials per place. Step 3 needs at most $16$ local uniformizers, $16^2$ valuations at the candidates, and $16$ valuations of pulled-back functions at $P$. The pulled-back functions are the expensive part, which is why Step 3 is used only when Steps 1--2 leave several candidates. The script reports, for each degree, how many images were decided by each step.

%======================================================================
\section{Checks}\label{sec:checks}
%======================================================================

The following are verified by \code{X1\_42\_orbits\_merge.py} before a degree is saved.
\begin{enumerate}[label=(C\arabic*)]
\item \emph{Group.} The loaded $\sigma_4$ has order exactly $12$ (in every job).
\item \emph{Consistent numbering.} All parallel jobs for a degree use the same list $\Pl_v$. Each job writes $\#\Pl_v$ and the checksum
\[
\mathrm{CHK}=\Bigl(\cdots\bigl((0\cdot B+h_1)\,B+h_2\bigr)\cdots\Bigr)\bmod M,\qquad h_i=\bigl(b_i\bmod p\bigr)+7\bigl(\kappa_i\bmod p\bigr),
\]
with $B=1000003$, $p=10^9+7$, $M=2^{61}-1$, and $b_i,\kappa_i$ the bucket and key of $P^{(v)}_i$. The merge requires all parts to agree.
\item \emph{Completeness.} Every index has exactly one image, computed by the job responsible for it, and no image is ambiguous.
\item \emph{Permutation.} The images form a permutation of $\{1,\dots,\#\Pl_v\}$. By Lemma~\ref{lem:basic}(c), this must hold, so it is a non-trivial check of the whole identification.
\item \emph{Group law.} All cycle lengths divide $12$ (Corollary~\ref{cor:cycles}).
\item \emph{Known case.} For $v=1$, the permutation equals $A_4$ in \eqref{eq:A4} or its inverse. This is a consistency check with \code{X1\_42\_check\_orbs.m} only: $A_4$ refers to the numbering of \code{Places(AFF,1)} in that session, and the inverse is accepted because replacing $\sigma$ by $\sigma^{-1}$ replaces $\pi_v$ by $\pi_v^{-1}$, with the same orbits.
\end{enumerate}
If any check fails for a degree, no file is written for that degree, and an old file of that degree is deleted, so a stale result can never be reused.

Two further checks are made later by \code{X1\_42\_rational\_labels.m}, independently of the session in which the orbit files were made:
\begin{enumerate}[label=(C\arabic*),resume]
\item \emph{Same generator and convention in all degrees.} For every place $Q$ of degree $2\le v\le5$ (recomputed in the current session), the pair $(b,\kappa)$ of $\sigma(Q)$ is computed from $\sigma^*t,\sigma^*w$ at $Q$ (Lemmas~\ref{lem:bucket} and~\ref{lem:key}) and compared with the pair of the image stored in \code{X1\_42\_orbits\_d<v>.txt}. Since $(b,\kappa)$ is unique in these degrees (Remark~\ref{rem:collisions}), agreement for all places proves that the stored permutations are exactly those of $\sigma_4$, with the same convention ($\sigma$, not $\sigma^{-1}$) as the degree-$1$ permutation~\eqref{eq:S}. This is recorded in the lines \code{V v count} of \code{X1\_42\_rational\_labels.txt}.
\item \emph{Rational places.} The permutation \eqref{eq:S} on the signature labels is computed directly from the pull-backs and checked to be a $12$-cycle.
\end{enumerate}

%======================================================================
\section{Parallel computation}\label{sec:merge}
%======================================================================

\subsection{Splitting}
The images of different places are independent, so the work splits trivially. With $m$ parts for degree $v$, job $k\in\{0,\dots,m-1\}$ computes $\pi_v(i)$ exactly for
\[
i\equiv k+1\pmod m.
\]
This round-robin split balances the load better than contiguous blocks, because Magma's order of $\Pl_v$ groups places by the place of $\F_5(t)$ below them, and bucket sizes vary.

Each job still computes the full list $\Pl_v$ and the buckets and keys of \emph{all} places, since any of them may be a candidate. It also loads the group and checks the order. These shared costs are small compared with the image computations for $v=4,5$.

\subsection{Launch and merge}
\code{run\_orbits.sh P1 P2 P3 P4 P5} starts $P_v$ jobs for degree $v$; the default is $1,1,1,4,12$, i.e.\ $19$ jobs, one per core:
\[
\code{magma -b degs:=}v\ \code{part:=}k\ \code{nparts:=}m\ \code{X1\_42\_orbits.m}
\]
Job $(v,k,m)$ writes \code{X1\_42\_orbits\_d<v>\_part<k>\_of\_<m>.txt}, which contains:
\begin{itemize}
\item a header with \code{DEG}, \code{GEN}, \code{PART}, \code{NPARTS}, \code{COUNT} and \code{CHECKSUM};
\item for part $0$ only, the place data lines \code{P i | c0 c1 ... | key};
\item its lines \code{I i j}, meaning $\pi_v(i)=j$;
\item a closing \code{END} line.
\end{itemize}
The log of each job is \code{log\_orbits\_d<v>\_part<k>\_of\_<m>.txt}. After all jobs finish, \code{X1\_42\_orbits\_merge.py} performs checks (C2)--(C6), splits $\pi_v$ into cycles, and writes the final files.

%======================================================================
\section{The saved data}
%======================================================================

For each degree $v$, the file \code{X1\_42\_orbits\_d<v>.txt} contains:
\begin{center}\small
\begin{tabular}{@{}lp{0.5\textwidth}@{}}
\toprule
line & meaning\\\midrule
\code{X1\_42\_ORBITS v1 DEG v GEN 4} & header, continued by\\
\quad\code{COUNT n ORBITS k CHECKSUM c} & \\
\code{P i | c0 c1 ... cn | key} & place $i$: coefficients of $\minpoly(t(P_i))$, constant term first (empty for $t=\infty$), and $\kappa(P_i)$\\
\code{I i j} & $\sigma_4(P_i)=P_j$\\
\code{O k s : i1 i2 ... is} & orbit $k$ of size $s$, as the cycle $i_1\mapsto i_2\mapsto\cdots$\\
\code{R k i} & representative of orbit $k$ (its smallest index)\\
\code{END n} & end marker\\\bottomrule
\end{tabular}
\end{center}
The indices $i$ are positions in Magma's list $\Pl_v$ in the session that made the file. For $2\le v\le5$ each index is tied to its session-independent data $(b,\kappa)$ by the \code{P} line, so the file can be read in any session (Section~\ref{sec:labels}). The \code{R} lines are not used by the search, which chooses its own representatives (Section~\ref{sec:search}). \code{X1\_42\_orbits\_summary.txt} lists, for each degree, the number of orbits and the distribution of orbit sizes.

\subsection{Reuse}
\code{X1\_42\_orbits\_load.m} provides three functions.
\begin{itemize}
\item \code{LoadOrbitFile(v)} reads the file into a record with fields \code{coeffs}, \code{keys}, \code{img}, \code{orbits}, \code{reps} and \code{orbitof}. It needs no curve computation.
\item \code{PlacesFromOrbitFile(D)} recomputes $\Pl_v$ and verifies, place by place, the stored polynomial and key and the checksum. If this succeeds, index $i$ in the file refers to \code{PL[i]} in the current session.
\item \code{FindPlace(coeffs, key, v)} recovers a single place without computing $\Pl_v$. It takes the zeros of $p(t)$, where $p$ is the stored polynomial (or the poles of $t$ if the data is empty), keeps those of degree $v$, and keeps the one with matching key. It succeeds when the stored data singles out the place, which is the case for every place of degree $2\le v\le5$.
\end{itemize}
\code{X1\_42\_orbits\_check.m} runs both reloading routes on every saved degree and prints PASS/FAIL lines. The auxiliary script \code{X1\_42\_orbit\_reps.m} writes \code{X1\_42\_orbit\_reps\_d<v>.txt} with an explicit description of one place in every orbit; it is not needed by the search.

%======================================================================
\section{Labels independent of the session}\label{sec:labels}
%======================================================================

Because Magma's order of $\Pl_v$ may change between sessions, the search refers to places by labels that depend only on the place.

\emph{Degrees $2\le v\le5$.} The label of a place $Q$ is the rank of $(b(Q),\kappa(Q))$ in sorted order among the places of degree $v$. This is well defined by Remark~\ref{rem:collisions}. \code{X1\_42\_split.py} translates the permutation $\pi_v$ of each orbit file into a permutation of these labels, and checks again that it is a permutation with all cycle lengths dividing $12$.

\emph{Degree $1$.} Ten of the twelve rational places share their pair $(b,\kappa)$, so they are labelled instead by \emph{signatures}. \code{X1\_42\_rational\_labels.m} fixes a list $g_1,\dots,g_n$ of functions: $t$, $w$, a local uniformizer at each rational place, and a basis of $\LL(3S)$, where $S$ is the sum of the rational places. It writes them to \code{X1\_42\_rational\_labels.txt} as explicit elements $\sum_{j=0}^{15}c_j(t)\,w^j$ (lines \code{F i data}), so that every later session reloads exactly the same functions. The signature of a rational place $P$ records, for each $g_i$, the valuation $v_P(g_i)$ and, if it is $0$, the value $g_i(P)\in\F_5$. The twelve signatures are distinct, and the label of $P$ is the rank of its signature in sorted order (lines \code{R k signature}). The permutation of $\sigma_4$ on these labels is computed from the signatures of the pulled-back functions $\sigma^*g_i$, using Lemma~\ref{lem:basic}(a),(b), and gives \eqref{eq:S} (line \code{S}). The line \code{Q1} records the label ($5$) of the zero of order $3$ of $w$, the place $Q_1$ of the manuscript; the pole of order $7$ of $w$, the place $Q_2$, has label $1$, and $\sigma_4^5(Q_2)=Q_1$.

%======================================================================
\section{Use in the divisor search}\label{sec:search}
%======================================================================

Let $S$ be the sum of the $12$ rational places. Since $G$ permutes the rational places, $\phi(S)=S$ for all $\phi\in G$, and for $\phi\in G$ the map $h\mapsto h\circ\phi^{-1}$ is a degree-preserving bijection $\LL(S+E)\to\LL(S+\phi(E))$ (Lemma~4.1 of the manuscript).

Every effective divisor $E$ of degree $9$ decomposes uniquely as $E=E_{\le5}+E_{\ge6}$, with $E_{\le5}$ supported on places of degree $\le5$ and $E_{\ge6}$ either $0$ or a single place of degree $6,\dots,9$ (since $2\cdot6>9$). By Theorem~6.2 of the manuscript, to prove $\gon_{\F_5}(X)\ge12$ it suffices to test, for one representative $E_{\le5}$ of every $G$-orbit of low-degree parts, all divisors $E_{\le5}+E_{\ge6}$, with $E_{\ge6}$ running over $0$ and all places of degree $6,\dots,9$. The action of $\sigma$ on places of degree $\ge6$ is never needed.

Write $E_{\le5}$ as the tuple, over the degrees $v\le5$ occurring in it in decreasing order, of the sorted lists of labels (Section~\ref{sec:labels}) of its places of degree $v$, each place listed as often as its multiplicity. Order such tuples lexicographically.

\begin{proposition}\label{prop:canon}
Keep $E_{\le5}$ if and only if its tuple is lexicographically smallest among the tuples of $\sigma^k(E_{\le5})$, $k=0,\dots,11$. Then every $G$-orbit of effective divisors supported on places of degree $\le5$ contains exactly one kept divisor.
\end{proposition}

\begin{proof}
$\sigma$ preserves degrees, so the tuples of $\sigma^k(E_{\le5})$ have the same shape, and are obtained by applying the stored permutations to the labels and re-sorting. Within an orbit the tuple determines the divisor, so the orbit has a unique lexicographically smallest tuple, and its divisor is the only kept element of the orbit.
\end{proof}

This is the rule defining $\mathcal{R}$ in Section~6.2 of the manuscript, implemented in \code{kept\_low} of \code{X1\_42\_split.py}. Since the permutations on the labels are stored, it needs no computation on the curve. As a check, \code{X1\_42\_split.py} compares the number of kept parts of each type with the number of $G$-orbits given by Burnside's lemma from the cycle types of $\sigma$ alone. The two agree for every partition of $9$ (Table~3 of the manuscript). In total $743\,231$ divisors are tested, compared with $2\,994\,534$ in the first search and $6\,370\,940$ without any reduction. The divisors with $E_{\ge6}\neq0$ are split among the jobs by the bucket of the high-degree place, as described in Section~6.3 of the manuscript.

%======================================================================
\section{Results}
%======================================================================

The orbits of $\sigma_4$ on the places of degree $v\le5$ (Appendix~A.2 of the manuscript) are:
\begin{center}
\begin{tabular}{@{}cccl@{}}
\toprule
$v$ & $\#\Pl_v$ & orbits & orbit sizes\\
\midrule
$1$ & $12$ & $1$ & $12$\\
$2$ & $6$ & $1$ & $6$\\
$3$ & $8$ & $2$ & $4,4$\\
$4$ & $183$ & $17$ & $3,\ 6^2,\ 12^{14}$\\
$5$ & $648$ & $54$ & $12^{54}$\\
\bottomrule
\end{tabular}
\end{center}
All orbit sizes divide $12$. In degree $4$ the group does not act freely: one orbit consists of places fixed by $\sigma^3$, and two of places fixed by $\sigma^6$.

\emph{Notes on Magma.} \code{LocalUniformizer} is used in Step~3, and the scripts check that it has positive valuation at the place. \code{UniformizingParameter} and \code{WeakApproximation([Q],[1])} are not available with these argument types in Magma V2.29-10. Loading the group takes about $2$\,s, and the order check about $87$\,s.

\end{document}
